\section{Admissible Spectral Orders}
\label{sec:orders}

The native order of the terminal residues is determined by traversal of the
factor tree.  Its apparent shuffle becomes a precise ordering problem once
locality is defined algebraically.

\begin{definition}[Subtree-contiguous order]
A total order on the leaves of a rooted binary factor tree is
subtree-contiguous if the leaves below every node form one interval of the
order.
\end{definition}

\begin{theorem}[Admissible orders]
An order is subtree-contiguous if and only if it is obtained by independently
choosing the left-right orientation at every internal node and then listing the
leaves recursively.
\end{theorem}
\begin{proof}
Every recursive orientation clearly produces intervals.  Conversely, the two
child leaf sets of the root are disjoint intervals whose union is the complete
order, so one precedes the other.  Restriction to each child is again
subtree-contiguous.  Induction gives an orientation at every internal node.
\end{proof}

This theorem explains the sorting experiments without referring to a storage
device.  Reordering siblings preserves the same local factor operations;
interleaving leaves from unrelated subtrees fragments at least one factor
interval.  The admissible freedom is exactly a word of sibling reversals.

\subsection{Frequency labels}

The label of a quadratic leaf is an unordered conjugate pair
$\{k,N-k\}$.  The binomial spine supplies a dyadic shell before the half-angle
word begins.  Write
\begin{equation}
 k=2^v u,\qquad u\text{ odd},\qquad0\le v\le L-2.
 \label{eq:dyadic-shell}
\end{equation}
The shell $v$ leaves the spine through
$z^{N/2^{v+1}}+1$.  Within that factor, start at $\alpha_0=\pi/2$ and apply
$\ell=L-v-2$ half-angle decisions
\begin{equation}
 \alpha\longmapsto {\alpha\over2},
 \qquad
 \alpha\longmapsto \pi-{\alpha\over2}.
 \label{eq:angle-address}
\end{equation}
After $\ell$ decisions, $k=N\alpha/(2\pi)$ has dyadic valuation $v$.  Choosing the
representative $1\le k<N/2$ and adjoining $0$ and $N/2$ defines a bijection from
native leaves to standard packed positions.

\begin{proposition}[Address bijection]
The shell index $v$ together with the recursive law
\eqref{eq:angle-address} visits every conjugate pair $\{k,N-k\}$ exactly once.
\end{proposition}
\begin{proof}
Every $1\le k<N/2$ has a unique dyadic valuation $v$.  The roots of
$z^{N/2^{v+1}}+1$ are precisely the frequencies $2^v u$ with $u$ odd.  By the
half-angle factorization, the two child root sets are disjoint and their union
is the parent root set.  The $\ell$ decisions terminate in every conjugate pair
of that shell once.  Unique valuation separates the shells.
\end{proof}

\begin{table}[t]
\caption{Native left-first factor-tree traversal for $N=16$.  A quadratic leaf
$q_k$ maps to packed coordinates $(2k-1,2k)$.}
\label{tab:n16-labels}
\centering
\begin{tabular}{@{}rrll@{}}
\toprule
Leaf & Shell $v$ & Factor & Frequency\\
\midrule
0 & -- & $z-1$ & $0$\\
1 & -- & $z+1$ & $8$\\
2 & 2 & $q_4$ & $4$\\
3 & 1 & $q_2$ & $2$\\
4 & 1 & $q_6$ & $6$\\
5 & 0 & $q_1$ & $1$\\
6 & 0 & $q_7$ & $7$\\
7 & 0 & $q_3$ & $3$\\
8 & 0 & $q_5$ & $5$\\
\bottomrule
\end{tabular}
\end{table}

The address algebra is therefore an exact boundary permutation.  It does not
alter the factorization, its inverse, or its metric.
